\subsection{Separate Chaining}
\begin{frame}{Separate Chaining의 구조}
\centering\begin{tikzpicture}[x=1.8cm,y=.72cm]
\foreach \i in {0,...,6}{\node[ht slot](b\i)at(0,-\i){\i};}
\node[chain node](a)at(2,-3){24};\node[chain node](c)at(4,-3){17};\node[chain node](d)at(6,-3){10};
\draw[chain edge](b3)--(a);\draw[chain edge](a)--(c);\draw[chain edge](c)--(d);
\node[chain node](e)at(2,-5){12};\node[chain node](f)at(4,-5){5};
\draw[chain edge](b5)--(e);\draw[chain edge](e)--(f);
\end{tikzpicture}
\begin{center}\small $h(k)=k\bmod7$, head insertion: bucket 3 chain은 $24\to17\to10$.\end{center}
\end{frame}
\begin{frame}{Chained Hash Operations}
\small\begin{algorithmic}[1]
\Procedure{ChainedSearch}{$T,key$}
 \For{each entry $e$ in $T[h(key)]$}
  \If{$e.key=key$}\State\Return $e$\EndIf
 \EndFor
 \State\Return NOT\_FOUND
\EndProcedure
\Statex
\Procedure{ChainedPut}{$T,key,value$}
 \State $e\gets\Call{ChainedSearch}{T,key}$
 \If{$e\ne$ NOT\_FOUND}\State $e.value\gets value$\Else
 \State insert $(key,value)$ at chain head; $n\gets n+1$
 \EndIf
\EndProcedure
\end{algorithmic}
\end{frame}
\begin{frame}{Chaining Animation: Insert}
\centering\begin{tikzpicture}[x=1.7cm,y=.8cm]
\foreach \i in {0,...,6}{\node[ht slot](b\i)at(0,-\i){\i};}
\only<1|handout:0>{\node[chain active]at(3,-3){10};\draw[chain edge](b3)--(1.8,-3);\node[callout]at(5,-1){$10\bmod7=3$};}
\only<2|handout:0>{\node[chain active](a)at(2,-3){17};\node[chain node](c)at(4,-3){10};\draw[chain edge](b3)--(a);\draw[chain edge](a)--(c);\node[callout]at(5,-1){collision; head insert 17};}
\only<3|handout:1>{\node[chain active](a)at(2,-3){24};\node[chain node](c)at(4,-3){17};\node[chain node](d)at(6,-3){10};\draw[chain edge](b3)--(a);\draw[chain edge](a)--(c);\draw[chain edge](c)--(d);\node[callout]at(5,-1){collision; head insert 24};}
\end{tikzpicture}
\end{frame}
\begin{frame}{Chaining Animation: Search와 Delete}
\centering\begin{tikzpicture}[x=1.9cm]
\node[ht slot](b)at(0,0){3};\node[chain node](a)at(2,0){24};\node[chain node](c)at(4,0){17};\node[chain node](d)at(6,0){10};
\draw[chain edge](b)--(a);\draw[chain edge](a)--(c);\draw[chain edge](c)--(d);
\only<1|handout:0>{\node[chain active]at(a){24};\node[callout]at(3,-1.2){SEARCH 17: compare 24};}
\only<2|handout:0>{\node[chain active]at(c){17};\node[callout]at(3,-1.2){17 found};}
\only<3|handout:0>{\node[chain active]at(c){17};\node[callout]at(3,-1.2){DELETE 17: unlink};}
\only<4|handout:1>{\draw[chain edge,orange](a)to[bend left=25](d);\node[chain node,fill=AlgoGray!10,text=AlgoGray]at(c){removed};\node[callout]at(3,-1.2){24 $\to$ 10};}
\end{tikzpicture}
\end{frame}
\begin{frame}{Chain Collection 선택}
\begin{itemize}
\item singly linked list: 단순, head insertion 가능
\item dynamic array/small vector: locality가 좋을 수 있음
\item sorted chain: search 조기 종료 가능하지만 insert 비용 증가
\item extreme collision에서 balanced tree를 고려할 수 있음
\end{itemize}
\begin{alertblock}{Map semantics}head insertion 자체는 $\Theta(1)$이지만 duplicate update를 지원하려면 먼저 equal key를 찾아야 한다.\end{alertblock}
\end{frame}
\begin{frame}{Chaining Analysis}
\[
\alpha=\frac{n}{m}=\text{average chain length}.
\]
\small\centering\begin{tabular}{lc}\toprule
operation (resize 제외)&simple uniform hashing 아래 expected\\\midrule
unsuccessful search&$\Theta(1+\alpha)$\\
successful search&$\Theta(1+\alpha)$\\
unchecked head insert&$\Theta(1)$\\
map put/update&$\Theta(1+\alpha)$\\\bottomrule
\end{tabular}
\httakeaway{$\alpha$를 bounded constant로 유지하면 resize 없는 연산은 expected $\Theta(1)$. 특정한 $M\approx N/10$이 보편 최적이라는 뜻은 아니다.}
\end{frame}
\begin{frame}{Chaining Worst Case}
\centering\begin{tikzpicture}[x=1.2cm]
\node[ht collision](b)at(0,0){bucket $j$};
\foreach \i in {1,...,6}{\node[chain node,minimum width=14mm](n\i)at({1.55*\i},0){$k_{\i}$};}
\draw[chain edge](b)--(n1);\foreach \i in {1,...,5}{\pgfmathtruncatemacro{\j}{\i+1}\draw[chain edge](n\i)--(n\j);}
\end{tikzpicture}
\[
\text{all keys in one bucket}\quad\Longrightarrow\quad \Theta(n)\text{ search}.
\]
\end{frame}
\begin{frame}{Separate Chaining Checkpoint}
\begin{enumerate}\item $m=7,n=12$이면 $\alpha$는?\item head insertion map put이 항상 $\Theta(1)$이 아닌 이유는?\item bucket 3의 search order는?\end{enumerate}
\begin{exampleblock}{답}$12/7$; duplicate/update search가 필요; $24,17,10$.\end{exampleblock}
\end{frame}
